\documentclass[11pt,a4paper]{ltjsarticle}
\usepackage[margin=24mm]{geometry}
\usepackage{amsmath,amssymb,amsthm,mathtools,booktabs}
\usepackage[haranoaji]{luatexja-preset}
\usepackage[unicode,hidelinks]{hyperref}
\newtheorem{theorem}{定理}[section]
\newtheorem{proposition}[theorem]{命題}
\newtheorem{lemma}[theorem]{補題}
\newtheorem{corollary}[theorem]{系}
\theoremstyle{definition}
\newtheorem{definition}[theorem]{定義}
\newtheorem{remark}[theorem]{注意}
\renewcommand{\proofname}{証明}
\newcommand{\ZZ}{\mathbb Z}
\newcommand{\BB}{\mathcal B}
\newcommand{\DD}{\mathcal D}
\newcommand{\CC}{\mathbb C}
\newcommand{\GL}{\operatorname{GL}}
\newcommand{\PGL}{\operatorname{PGL}}
\newcommand{\Gr}{\operatorname{Gr}}
\newcommand{\Aut}{\operatorname{Aut}}
\newcommand{\Per}{\operatorname{Per}}
\newcommand{\qgr}{\operatorname{qgr}}
\title{四次元二次AS正則\texorpdfstring{$\ZZ$}{Z}-代数の分類\\
\large 4周期表示・1周期構造の忘却・全分類への有限atlas}
\author{}
\date{2026年11月12日\quad 研究原稿 v3（完全分類は未完成）}
\begin{document}
\maketitle
\begin{abstract}
本稿は標準$(1,4,6,4,1)$型のquadratic Koszul AS正則$\ZZ$-代数の
完全分類を目標とする研究原稿である。
4周期同型を\emph{選んだ}二次$\ZZ$-代数が四つの4次元生成空間と
四つの6次元関係空間によって記述されることを証明し、その群作用を同定する。
この有限atlasは、選んだ4周期同型を\emph{忘れた}対象のモジュライではない。
既存のgraded AS4分類の結果を1周期的部分の研究基盤として利用し、
既に構成された最小周期2、4の反例を全体問題に位置づける。
AS正則性の全条件、周期構造忘却のファイバー、全閉値の同型条件、
全対象の被覆および既約成分の分類は未完成として明記する。
\end{abstract}
\tableofcontents

\section{Introduction：対象と分類問題}
本稿を通じて$k$は標数零の代数閉体とする。
\emph{正連結な$\ZZ$-代数}$C=(C_{ij})_{i,j\in\ZZ}$とは
\[
 C_{ii}=k,\quad C_{ij}=0\ (i>j),\quad
 C_{ij}C_{j\ell}\subseteq C_{i\ell}
\]
を満たす局所有限な代数である。$C(s)_{ij}=C_{i+s,j+s}$とする。
\[
 \Per(C)=\{s\in\ZZ:C\xrightarrow{\sim}C(s)
 \text{という添字保持同型が存在する}\}
\]
を周期群と呼ぶ。
通常のconnected graded algebra $A=\bigoplus_{n\ge0}A_n$からの
付随$\ZZ$-代数$\check A$は$\check A_{ij}=A_{j-i}$である。
\emph{1周期同型の存在}と、その\emph{具体的な選択}は区別する。

本研究の中心対象を\emph{標準14641型}と呼ぶ。
これは次数一生成、二次関係のみを持ち、
\[
 \dim_k C_{i,i+n}={n+3\choose3}\quad(n\ge0)
\]
であり、各右単純加群$S_i$に有限最小線型解消
\[
 0\longrightarrow P_{i+4}
 \longrightarrow P_{i+3}^{\,4}
 \longrightarrow P_{i+2}^{\,6}
 \longrightarrow P_{i+1}^{\,4}
 \longrightarrow P_i\longrightarrow S_i\longrightarrow0
\]
が存在し、左側の対応解消とAS双対性を満たすものをいう。
ここで$P_i=e_iC$である。
これはKoszul型であり、大域次元4を意図する。
数値だけを仮定する弱い意味のquadratic AS正則代数との違いは保持する。
domain、noetherian、PBW、Calabi--Yau性を全域に仮定しない。

本研究では、添字保持$\ZZ$-代数同型を同値関係とする。
graded algebra同型を分類する\texttt{as4}の成果を参照するが、
その分類は現在なお未完成である
（\href{https://github.com/uedakazushi/as4}{\texttt{uedakazushi/as4}}の現行研究方針）。
SierraやZhangのtwisting systemによる同一視は、
graded同型より粗い$\ZZ$-代数同型を生む
\cite{Zhang,Sierra,MoriNyman}。

\subsection*{既知の出発点と本稿の新しい位置づけ}
標準14641型の4周期性はKoszul--Frobenius双対性から従い、
最小正周期は1、2、4のいずれかである
\cite{BP,MoriNyman}。
\texttt{as4}の
\href{https://github.com/uedakazushi/as4/blob/main/research/z_periodicity/papers/as4_z_periodicity_ja.tex}
{周期性研究原稿}には、最小周期2の平坦PBW族、
最小周期4の明示例、全閉値でのKoszul/AS正則性の証明、
有限恒等式の検証コードが記されている。
この原稿の独立査読・結果の新規性の最終判定は未了である。
特に、1周期的な特殊ファイバーと非1周期的な一般ファイバーを
結ぶ平坦族があるので、周期1/2/4を既約成分のラベルだとみなしてはならない。
本稿はそれらの証明を重複せず、**周期の選択を忘れる際の問題**に焦点を当てる。

既存の\texttt{as4}のgraded結果をすべて分類し終えても、
最小周期2、4の対象はそこから復元できない。
一方で、それらが必ず新しい既約成分を与えるとも限らない。
Bondal--Polishchukのhelix・quantum determinantの方法\cite{BP}と、
正則$\ZZ$-空間の圏的構成\cite{MoriNyman}を使うのが自然である。

\section{周期群と4周期の有限表示}
\begin{lemma}\label{lem:period-subgroup}
$\Per(C)\subset\ZZ$は部分群である。
特に$C$が4周期的なら、その最小正周期は$1,2,4$のいずれかである。
\end{lemma}
\begin{proof}
$s$周期同型と$t$周期同型は、後者を$s$だけ平行移動させて合成できるから
$s+t$は周期である。逆同型を$-s$だけ移動させると$-s$も周期となる。
$0$は自明に周期である。よって$\Per(C)=p\ZZ$。
$4\in p\ZZ$なら$p\mid4$となる。
\end{proof}

以下、\emph{4周期同型付き}対象$(C,\tau)$とは
\[
 \tau_{ij}:C_{ij}\xrightarrow{\sim}C_{i+4,j+4}
\]
が積を保つものとする。対象の同型は$\tau$と可換する添字保持
$\ZZ$-代数同型である。
$\tau$の存在のみを要求する裸の$C$とは別である。

四つの4次元$k$-空間$V_0,V_1,V_2,V_3$を固定し、添字を法4で読む。
\[
 \mathfrak P=\prod_{i=0}^{3}\Gr(6,V_i\otimes V_{i+1}),
 \qquad G=\prod_{i=0}^3\GL(V_i)
\]
とおく。$\Gr(6,16)$の次元は60なので、$\dim\mathfrak P=240$、
$\dim G=64$である。
ただし、この差176はAS正則モジュライの次元ではない。
AS条件をまだ課しておらず、安定化群もあるためである。

\begin{theorem}[4周期表示の有限四窓atlas]\label{thm:four-window}
各$C_{i,i+1}$が4次元で各二次関係空間が6次元の
\emph{次数一生成・二次}正連結$\ZZ$-代数$(C,\tau)$で、
4周期同型$\tau$を指定したものの同型群oidは、
$\mathfrak P$の$k$-点に対する$G$の作用群oidと同値である。
この定理ではAS正則性やKoszul性を仮定しない。
それらの条件を満たす全部分群oidにも同じ対応が制限される。
\end{theorem}
\begin{proof}
$(C,\tau)$が与えられたとき、$C_{i,i+1}$（$0\le i<4$）の基底を選び
$V_i$と同一視する。$\tau$の反復により全整数$i$へ基底を延長する。
$C$の二次関係空間
\[
 R_i=\ker\left(C_{i,i+1}\otimes C_{i+1,i+2}
 \longrightarrow C_{i,i+2}\right)
\]
は$\tau$と両立するので$R_{i+4}=R_i$となり、
$\dim R_i=6$である。従って$\mathfrak P(k)$の点を定める。

逆に$(R_0,R_1,R_2,R_3)$が与えられれば、整数列上、
頂点$i$から$i+1$へ$V_{i\bmod4}$の四つの矢印を置く。
自由な経路$\ZZ$-代数の中で、
各$i$の$R_{i\bmod4}\subset V_i\otimes V_{i+1}$から生成される
二次両側イデアルを商に取る。
この商は正連結、次数一生成、二次で、
二次成分の核はちょうど与えた$R_i$である。
係数が4ごとに一致するので、添字を4移す厳密な同型$\tau$を持つ。

選んだ基底を変えると$G$の自然な作用になる。
逆向きに$\tau$と可換する同型は、全ての次数一生成元上で
四つの可逆写像$g_i:V_i\to V'_i$から定まり、
二次関係を保つことは
$(g_i\otimes g_{i+1})(R_i)=R'_i$に同値である。
二次生成性からこの条件は全積を保つ同型の十分条件でもある。
同型の合成も同じなので群oid同値が従う。
\end{proof}

\begin{remark}[分類上の重要な違い]\label{rem:forget4}
定理\ref{thm:four-window}は、
4周期同型$\tau$を選んだ$(C,\tau)$の分類であり、
裸の$C$に対する全単射を与えない。
同じ$C$の上に相異なる4周期同型が存在し得る。
周期性の存在とその選択の違いを処理しない限り、
$[\mathfrak P/G]$を求めても全$\ZZ$-代数の分類は完成しない。
また、$\mathfrak P$の大部分がKoszul AS正則であるとは主張していない。
\end{remark}

\begin{proposition}[次数一の関係式による周期の判定]\label{prop:period-relations}
二次$\ZZ$-代数$C$について、整数$s$が周期であることは、
全$i\in\ZZ$に対する可逆線型写像
$g_i:V_i=C_{i,i+1}\to V_{i+s}=C_{i+s,i+s+1}$が存在して
\[
 (g_i\otimes g_{i+1})R_i=R_{i+s}
\]
を満たすことと同値である。
\end{proposition}
\begin{proof}
$\ZZ$-代数同型を次数一に制限すれば必要性は直ちに従う。
逆に与えられた可逆写像は自由な経路$\ZZ$-代数の同型を定め、
各二次関係生成空間を対応する関係生成空間へ送る。
従って二次イデアル全体を保ち、商に同型を誘導する。
\end{proof}

\begin{remark}
命題\ref{prop:period-relations}の写像$g_i$は、
あらかじめ選んだ4周期$\tau$と\emph{可換する必要がない}。
最初から$g_{i+4}=g_i$という条件だけで1周期性を調べると、
選んだ$\tau$と可換しない周期同型を見落とす。
この点が未標識moduliの分類で最初に解決すべき障害となる。
\end{remark}


\section{Koszul--Frobenius標準4周期構造と忘却}\label{sec:canonical-four}
以下の結果は、4周期同型を任意に選ぶことの余分な自由度を
\emph{幾何学的点の同型類}の水準で取り除く。
モジュライstackの同値や族上の大域的な標準化は、
この段階では主張しない。

$C$を標準14641型とし、前向き二次Koszul双対を
$E=C^!$と書く。ここでは双対の矢印を
$E_{i,i+1}=C_{i,i+1}^{*}$とする規約を採用する。
左右AS双対性とKoszul性により、各$i$に
\[
 E_{i,i+4}\simeq k,\quad
 \dim E_{i,i+p}={4\choose p}\ (0\le p\le4),
 \quad E_{i,i+p}=0\ (p>4)
\]
があり、任意の$p$について積による相補次数の対合が完全である
\cite{BP}。

各$i$に非零な最高次数のtrace
$\lambda_i:E_{i,i+4}\xrightarrow{\sim}k$を選ぶ。
相補次数の完全性から、一意的な線型同型
\[
 \nu^\lambda_{ij}:E_{ij}\xrightarrow{\sim}E_{i+4,j+4}
\]
が
\begin{equation}\label{eq:nakayama-pair}
 \lambda_i(ab)=\lambda_j\bigl(b\,\nu^\lambda_{ij}(a)\bigr),
 \quad a\in E_{ij},\ b\in E_{j,i+4}
\end{equation}
で定まる。この$\nu^\lambda$は積を保つ$\mathbb Z$-代数同型である。
実際、$a\in E_{ij}$、$c\in E_{j\ell}$、$b\in E_{\ell,i+4}$に対して
\[
 \lambda_i((ac)b)
 =\lambda_j((cb)\nu^\lambda(a))
 =\lambda_\ell(b\,\nu^\lambda(a)\nu^\lambda(c))
\]
なので、完全対合から
$\nu^\lambda(ac)=\nu^\lambda(a)\nu^\lambda(c)$となる。
次数一部分を逆双対化し、二次関係が保たれることから、
$C$にも4周期同型
\[
 \tau^\lambda:C\xrightarrow{\sim}C(4)
\]
が定まる。これを\emph{Frobenius標準4周期同型}と呼ぶ。

\begin{theorem}[標準4周期同型の選択独立性]\label{thm:canonical-four}
$C$を標準14641型のKoszul AS正則$\mathbb Z$-代数とする。
任意の二つの非零trace族$\lambda=(\lambda_i)$と
$\lambda'=(\lambda'_i)$から得られる
$\tau^\lambda$と$\tau^{\lambda'}$は、$C$の
添字保持自己同型$h$により
\[
 \tau^{\lambda'}=h(4)\,\tau^\lambda\,h^{-1}
\]
の形で共役になる。
さらに、$C$と$C'$が$\mathbb Z$-代数として同型なら、
それぞれのFrobenius標準4周期同型付きの組も同型である。
従って\emph{幾何学的点の同型類}に限れば、
裸の標準14641型$\mathbb Z$-代数の分類は、
Frobenius標準4周期同型を選んだ組の分類に一致する。
\end{theorem}
\begin{proof}
$\lambda'_i=t_i\lambda_i$（$t_i\in k^\times$）と置く。
式\eqref{eq:nakayama-pair}を双方で比較すると
\begin{equation}\label{eq:trace-change}
 \nu^{\lambda'}_{ij}=(t_i/t_j)\nu^\lambda_{ij}
\end{equation}
を得る。
各整数$i$にスカラー$s_i\in k^\times$を、
\[
 s_{i+4}/s_i=t_i
\]
となるように選べる。これは四つの剰余類でそれぞれ
$s_i$を一つ選び、正負両方向へ漸化式で延ばせばよい。
頂点スカラーによる自己同型
\[
 h^E_{ij}:E_{ij}\longrightarrow E_{ij},
 \qquad a\longmapsto(s_i/s_j)a
\]
は積を保つ。従って
\[
 h^E(4)\nu^\lambda(h^E)^{-1}(a)
 =\frac{s_{i+4}s_j}{s_{j+4}s_i}\nu^\lambda(a)
 =\frac{t_i}{t_j}\nu^\lambda(a)
 =\nu^{\lambda'}(a).
\]
二次双対により、$C$上の二つの$\tau$も
対応する頂点スカラー自己同型によって共役となる。
（双対化でスカラーの逆が現れるが、共役の向きは保たれる。）

次に$f:C\xrightarrow{\sim}C'$が与えられれば、
二次双対の同型を通じて$\lambda$を$E'=(C')^!$に移送できる。
式\eqref{eq:nakayama-pair}の特徴付けにより、
移送したtraceから得られる4周期同型と$\tau^\lambda$は
$f$と両立する。任意の$\lambda'$を指定した場合も
既に示した頂点スカラー共役で補正すればよい。
よって忘却はFrobenius標準の組の\emph{同型類の集合}上で
全単射である。これを群oidそのものの完全忠実性とはしていない。
\end{proof}

\begin{proposition}[有限四窓でのFrobenius標準化条件]\label{prop:canonical-matrix}
定理\ref{thm:four-window}の厳密4周期表示で、二次双対$E$が
上述の局所Frobenius条件を満たすとする。
4つの非零trace$\lambda_i:E_{i,i+4}\to k$
（$i\in\mathbb Z/4$）を選び、厳密4移動によって延長する。
$E_{i,i+1}$と$E_{i+1,i+4}$の基底を選び、
$E_{i+4,i+5}=E_{i,i+1}$を厳密移動で同定する。
\[
 P_i=(\lambda_i(u_a v_b))_{a,b},\qquad
 Q_i=(\lambda_{i+1}(v_b u'_c))_{b,c}
\]
を、$u_a\in E_{i,i+1}$、$v_b\in E_{i+1,i+4}$、
$u'_c\in E_{i+4,i+5}$の対応する基底により定める。
このとき両行列は4次正則行列で、
\[
 N_i:=Q_i^{-1}P_i^{\mathsf T}
\]
は$\nu^\lambda_{i,i+1}$を厳密4移動で同定した行列である。
従って選んだ厳密4周期同型がFrobenius標準になることは、
\emph{四つのtraceを適切に選んで}
\begin{equation}\label{eq:canonical-matrix-test}
 P_i^{\mathsf T}=Q_i\quad(i=0,1,2,3)
\end{equation}
を満たせることと同値である。
\end{proposition}
\begin{proof}
$N_i$の列を$\nu^\lambda(u_a)$の係数とすると、
式\eqref{eq:nakayama-pair}は各$b,a$について
$P_{i,ab}=\sum_c Q_{i,bc}(N_i)_{ca}$であり、
$P_i^{\mathsf T}=Q_iN_i$となる。
完全対合から$Q_i$は可逆である。
従って$N_i=I$は\eqref{eq:canonical-matrix-test}と同値。
$E$は次数一生成なので、$\nu^\lambda$が次数一で
厳密4移動に一致すれば、全積でも一致する。
その逆双対の$C$の4周期同型についても同様である。
\end{proof}

\begin{proposition}[有限次数で検査できるFrobenius条件]\label{prop:finite-frobenius}
任意の$\mathfrak P$の四窓データから構成する二次双対$E$について、
各始点$i\in\mathbb Z/4$で、
\[
 \dim(E_{i,i+p})=(1,4,6,4,1)_p\quad(0\le p\le4),
 \qquad E_{i,i+5}=0
\]
を検査し、各$0\le p\le4$の積
\[
 E_{i,i+p}\otimes E_{i+p,i+4}\longrightarrow E_{i,i+4}
\]
が完全対合であることを確認すれば、
$E$の局所Frobenius性は全次数で成立する。
これらは四窓の関係係数に対する\emph{有限回の線型代数}
（商空間の次元・積行列の階数）で判定できる。
ただしこの有限検査\emph{だけ}から元の二次代数$C$のKoszul性、
AS正則性は結論できない。
\end{proposition}
\begin{proof}
$E$は次数一生成であるため、各始点で長さ5の全語が零なら、
長さ6以上の語も先頭5語を用いて零となる。
有限次元の相補次数対合の完全性は仮定した積行列の
行列式の非消滅と同値である。
必要な関係空間、低次成分、積行列は、いずれも長さ5以下の
テンソル積空間から二次関係による部分空間を割るだけで
構成でき、4周期性により始点は4通りで足りる。
Koszul複体の正確性は別途要求されるので、
次数情報とFrobenius性だけをAS性へすり替えてはいない。
\end{proof}

\begin{corollary}[幾何学的点の標準化された分類問題]\label{cor:canonical-slice}
標準14641型の全$\mathbb Z$-代数の幾何学的同型類は、
$\mathfrak P$のうち\emph{Koszul AS正則かつ}
\eqref{eq:canonical-matrix-test}を満たすtraceの選択がある
四窓データの$G$-軌道と全単射になる。
ここで$G=\prod_{i=0}^{3}\GL(V_i)$は定理\ref{thm:four-window}の群である。
この軌道の幾何学的分類、AS locusの有限な必要十分条件、
相対族のstack構造はまだ未完成である。
\end{corollary}
\begin{proof}
任意の$C$には局所Frobenius traceから構成した4周期同型があるので
定理\ref{thm:four-window}で厳密4周期表示を選べる。
その表示は命題\ref{prop:canonical-matrix}を満たす。
定理\ref{thm:canonical-four}によって、
同じ裸の$C$から生じる二つのFrobenius標準表示は、
4周期同型を保つ同型で比較できる。
定理\ref{thm:four-window}の同型判定から、
それらは$G$の同一軌道になる。逆方向は忘却で得られる。
\end{proof}

\begin{theorem}[裸の周期性の有限四窓判定]\label{thm:finite-period-test}
$C$を標準14641型とし、Frobenius標準4周期同型を選んで
定理\ref{thm:four-window}の厳密4周期データ
$(R_0,R_1,R_2,R_3)$に直す。
任意の整数$p$に対して、
\[
 C\simeq C(p)
\]
となることは、四つの可逆線型写像
$g_i:V_i\xrightarrow{\sim}V_{i+p}$（$i\in\ZZ/4$）が存在して
\begin{equation}\label{eq:finite-period-test}
 (g_i\otimes g_{i+1})R_i=R_{i+p}
 \qquad(i\in\ZZ/4)
\end{equation}
を満たすことと同値である。添字は4を法として読む。
従って最小周期1、2、4の判定は、
Frobenius標準四窓上の\emph{有限個の多項式方程式と行列式非零条件}
による可解性判定に還元される。
\end{theorem}
\begin{proof}
\eqref{eq:finite-period-test}があれば$g_{i+4}=g_i$と延長し、
命題\ref{prop:period-relations}により
$C\xrightarrow{\sim}C(p)$を得る。

逆に$f:C\xrightarrow{\sim}C(p)$があるとする。
$\tau:C\to C(4)$を選んだFrobenius標準同型とすると、
そのシフト$\tau(p):C(p)\to C(p+4)$も
$C(p)$のFrobenius標準同型である。
$f$を使って$\tau$を$C(p)$へ移送して得られる4周期同型
$f(4)\tau f^{-1}$もFrobenius標準である。
定理\ref{thm:canonical-four}の証明で示した
\emph{頂点スカラーによる共役}を使えば、
$C(p)$の添字保持自己同型$h$が存在して
\[
 \tau(p)=h(4)\,[f(4)\tau f^{-1}]\,h^{-1}
\]
となる。$f'=hf$と置くと
\[
 \tau(p)\,f'=f'(4)\,\tau
\]
である。従って$f'$の次数一部分は
Frobenius標準4周期座標で$g_{i+4}=g_i$を満たす。
$f'$が二次関係を保つことから
\eqref{eq:finite-period-test}が従う。
一般の4周期同型を選んだ場合には、この補正が可能とは
限らないので、Frobenius標準という仮定が不可欠である。
\end{proof}

\begin{corollary}[回転軌道による周期層の記述]\label{cor:period-orbits}
Frobenius標準の四窓データに対して、
四つの関係空間の添字を一つ移す回転を$r$と書く。
このとき、裸の$\ZZ$-代数$C$の最小正周期は
\[
 \begin{array}{c|l}
 1& r(R)\text{が}R\text{と}G\text{-同一軌道},\\
 2& r^2(R)\text{が}R\text{と同一軌道で、}r(R)\text{は異軌道},\\
 4& r^2(R)\text{が}R\text{と異軌道}
 \end{array}
\]
で判定できる。
これらは同型類についての判定であって、
周期種別が既約成分を切り分けるとは主張しない。
\end{corollary}
\begin{proof}
定理\ref{thm:finite-period-test}の有限行列式条件が、
$R$と回転した$R$が$G$-同一軌道にあることと同値である。
補題\ref{lem:period-subgroup}によって最小周期は
$1,2,4$のいずれかである。
\end{proof}

\begin{remark}[有限atlasのスカラーkernel]
$G=\prod_{i=0}^{3}\GL(V_i)$の中で、
各$V_i$にスカラーとして作用する$(k^\times)^4$は
$\mathfrak P$の全ての点を固定する。
実際、$R_i\subset V_i\otimes V_{i+1}$には
$(s_i s_{i+1})$倍として作用するからである。
従って作用の有効部分は$\prod_i\PGL(V_i)$を通る。
ただしこの四次元トーラスは\emph{同型群oidの自己同型}
として残るため、stackを作る際に自動的に取り去ってはならない。
$\dim\mathfrak P=240$と有効群の次元60の差180を
AS正則モジュライの次元と呼ぶのも誤りである。
\end{remark}

\section{1周期的部分とgraded AS4の再利用}
\begin{proposition}\label{prop:graded-period}
二次正連結$\ZZ$-代数$C$は、
$C\simeq C(1)$であることと、ある次数一生成二次のconnected
graded algebra $A$の付随代数$\check A$と同型であることが同値である。
1周期同型を\emph{選んだ組}$(C,\sigma)$は、
graded代数の同型類と同値な分類問題を与える。
\end{proposition}
\begin{proof}
$C=\check A$のとき、添字を一つ移す恒等同型が存在する。
逆に$\sigma:C\to C(1)$を選び、$A_n=C_{0n}$として
$a\in A_m,b\in A_n$の積を
$a*b=a\,\sigma^m_{0n}(b)$と定める。
$\sigma$が積を保つことから結合律が従う。
$A_0=k$で、次数一生成・二次関係も引き継がれる。
$\sigma$を使い添字ごとの成分を同定すると
$\check A\simeq C$が得られる。
これらの構成を同型に適用すれば組とgraded代数の同値を得る。
\end{proof}

従ってgraded側である\texttt{as4}の定理を利用するには
既知族を$\ZZ$-代数へ写す際の\emph{忘却}を明示しなければならない。
gradedの同型条件を$\ZZ$-同型条件とみなしてはならない。
また、Zhang twistによる同型はperiodic presentationにより
同じ裸の$C$を与え得る。
Nakayama自己同型やCY性は選んだgraded表示に依存するため、
Z側の不変量としてそのまま採用しない。

\begin{proposition}[整合する周期同型と四窓の対角表示]\label{prop:compatible-periodic-windows}
$(C,\tau_4)$を4周期同型付き二次$\mathbb Z$-代数とし，
$p=1$または$2$とする。
$\sigma_p:C\xrightarrow{\sim}C(p)$が存在し，
その$4/p$回の合成が$\tau_4$に等しいとする。
このとき四窓の基底を適切に選ぶと，
$V_i=k^4$の同定の下で
\[
 R_i=R_{i+p}\quad(i\in\mathbb Z/4)
\]
が成立する。
逆に四窓の座標でこの等式が成立するならば，
指定した4周期同型$\tau_4$を$(4/p)$乗として与える
$p$周期同型が存在する。
\end{proposition}
\begin{proof}
$p$個の剰余類代表$i=0,\dots,p-1$の次数一空間に基底を選ぶ。
それらの基底を$\sigma_p$で整数全体に運ぶと，
$\sigma_p$は各次数一空間の間で同一の座標を持つ。
その$(4/p)$乗が$\tau_4$であるため，
この基底は定理\ref{thm:four-window}で選んだ4周期座標とも両立する。
$\sigma_p$が二次関係を保つので，$R_i=R_{i+p}$である。

逆に四窓の座標で$R_i=R_{i+p}$ならば，
頂点$i$から$i+p$への同一座標写像は二次関係空間を保つ。
命題\ref{prop:period-relations}によって
$\sigma_p:C\xrightarrow{\sim}C(p)$を定める。
その$(4/p)$乗は座標上の厳密な4移動，すなわち選んだ$\tau_4$である。
\end{proof}

\begin{remark}[1周期性と対角軌跡の違い]
$p=1$では四つの関係空間が同一になる\emph{完全対角}，
$p=2$では$(R_0,R_1,R_0,R_1)$型の\emph{二周期対角}が現れる。
しかしこれは$\tau_4=\sigma_p^{4/p}$という\emph{整合する選択}の結果である。
ある$C$が1周期的または2周期的であっても，
\emph{任意に選ばれた}$\tau_4$がその形の根を持つとは限らない。
従って完全対角軌跡の$G$-軌道を取るだけで，裸の$C$の
1周期的部分全体を分類したとは言えない。
\end{remark}

\section{既存の非1周期例と今後の完全分類}
既存研究\texttt{as4/research/z\_periodicity}では、
$\varepsilon_i=(-1)^i$として各位置の四生成元$X,Y,Z,W$に
\begin{align*}
 XY+YX&=Z^2+W^2,& XZ+ZX&=\varepsilon_i\lambda WY,\\
 XW+WX&=Z^2,&YZ+ZY&=W^2,\\
 YW+WY&=0,&ZW+WZ&=0
\end{align*}
を課した2周期PBW族が調べられている。
全ての閉パラメータで14641型・Koszul・左右AS正則であり、
$\lambda=0$は1周期、ある非空Zariski開集合では最小周期2である。
この結果は\texttt{as4}の原稿・有限検証コードを参照しており、
本稿が再び検証を実行したとは主張しない。
また別の例で最小周期4が実現される。
これらの例が全$\ZZ$-moduliのどの既約成分に属するか、
全特殊値の族内外同型がどうなるかは未解決である。

今後の分類では、(i) $\mathfrak P$に課すKoszul/AS条件、
(ii) $(C,\tau_4)\mapsto C$の同型群oidのファイバー、
(iii) 1/2/4の周期構造とgraded側の忘却、
(iv) 全特殊値の必要十分条件、
(v) 全対象被覆とモジュライ成分を順に調べる。
$\mathbb Z$-代数の最小周期が4を割ることだけでは
全てを分類したことにはならない。

\begin{thebibliography}{9}
\bibitem{BP}
A.~I.~Bondal and A.~E.~Polishchuk,
\emph{Homological properties of associative algebras: the method of helices},
Russian Acad. Sci. Izv. Math. \textbf{42} (1994), no.~2, 219--260.
\href{https://doi.org/10.1070/IM1994v042n02ABEH001536}{doi:10.1070/IM1994v042n02ABEH001536}.
\bibitem{MoriNyman}
I.~Mori and A.~Nyman,
\emph{A categorical characterization of quantum projective $\mathbb Z$-spaces},
Applied Categorical Structures \textbf{33} (2025).
\href{https://doi.org/10.1007/s10485-025-09806-2}{doi:10.1007/s10485-025-09806-2};
\href{https://arxiv.org/abs/2307.15253}{arXiv:2307.15253}.
\bibitem{Zhang}
J.~J.~Zhang, \emph{Twisted graded algebras and equivalences of graded categories},
Proc. Lond. Math. Soc. \textbf{72} (1996), 281--311.
\bibitem{Sierra}
S.~J.~Sierra, \emph{G-algebras, twistings, and equivalences of graded categories},
Algebras and Representation Theory \textbf{14} (2011), no.~2, 377--390.
\href{https://doi.org/10.1007/s10468-009-9193-y}{doi:10.1007/s10468-009-9193-y};
\href{https://arxiv.org/abs/math/0608791}{arXiv:math/0608791}.
\bibitem{as4period}
AS4周期性研究原稿, 
\emph{四次元二次AS正則$\mathbb Z$-代数の周期性：最小周期1,2,4と非1周期的な平坦族},
2026年11月12日・第2版,
\url{https://github.com/uedakazushi/as4/blob/main/research/z_periodicity/papers/as4_z_periodicity_ja.tex}.
\end{thebibliography}
\end{document}
